\documentclass[11pt]{article}
\usepackage[T1]{fontenc}
\usepackage{lmodern,microtype,amsmath,amssymb,amsthm,mathtools,booktabs,xcolor,fancyhdr}
\usepackage[letterpaper,margin=1.0in,headheight=14pt]{geometry}
\usepackage[colorlinks=true,urlcolor=blue!55!black,linkcolor=black,citecolor=blue!55!black]{hyperref}
\definecolor{ink}{HTML}{193A4B}
\pagestyle{fancy}\fancyhf{}\lhead{\small\sffamily Near the peak of a prime sum}\rhead{\small\sffamily Dirichlet characters}\cfoot{\thepage}
\newtheorem{theorem}{Theorem}\newtheorem{lemma}[theorem]{Lemma}\newtheorem{corollary}[theorem]{Corollary}
\newcommand{\eps}{\varepsilon}\newcommand{\rd}{\operatorname{rd}}\newcommand{\EE}{\mathbb E}\newcommand{\PP}{\mathbb P}
\newcommand{\dd}{\,\mathrm d}\newcommand{\meas}{\operatorname{meas}}
\newcommand{\Q}{\mathbb Q}\newcommand{\R}{\mathbb R}
\setlength{\emergencystretch}{2em}
\begin{document}
\thispagestyle{plain}
{\small\sffamily\color{ink} ANALYTIC NUMBER THEORY\hfill 26 NOVEMBER 2026}\par\vspace{10pt}
{\huge\bfseries\color{ink} Near the peak of a prime sum}\par\vspace{4pt}
{\large Logarithmic densities, rational relations among zeros, and the cyclotomic conductor}\par\medskip
\noindent\textbf{Abstract.}
A smoothed sum of prime powers in the class $1\bmod q$ returns to a limiting upper endpoint. We bound the logarithmic frequency of these returns using Dirichlet GRH alone. A finite-interval comparison works even when zero ordinates coincide or satisfy rational relations. If the loss of rational rank through height $T$ is $o(T/\log T)$, we determine the logarithm of the return frequency through its first conductor-dependent term. The first coefficient determines $\varphi(q)$, and the next arithmetic coefficient determines the root discriminant of $\Q(\zeta_q)$. Together they recover the cyclotomic modulus up to the unavoidable equivalence between an odd modulus and its double.

\section{The observable and its two density laws}
Fix $q\ge1$ and put $n=\varphi(q)$, with $n=1$ for $q=1$. Write $\Lambda$ for the von Mangoldt function. Form the real function
\[
 S_q(x)=\sum_{\substack{m\le x\\m\equiv1\bmod q}}\Lambda(m)(1/m-1/x),\quad
 R_q(x)=\sqrt{x}\bigl(S_q(x)-n^{-1}\log x-c_q\bigr).
\]
Here $c_q=\lim_{x\to\infty}(S_q(x)-n^{-1}\log x)$; Section~\ref{sec:source} gives a formula. Put $M_q=\limsup_{x\to\infty}R_q(x)$. Let $P_q^-(\eps)$ and $P_q^+(\eps)$ be the lower and upper limits, as $Y\to\infty$, of
\[
 \frac1Y\int_1^{e^Y}\mathbf1_{\{R_q(x)>M_q-\eps\}}\frac{\dd x}{x}.
\]
All limits in $\eps$ below are taken after these time limits. All logarithms are natural and $q$ is fixed.

The constants use the modified Bessel function $I_0$. Set
\begin{align*}
 g(u)&=u-\log I_0(u),&
 A&=\int_0^\infty g(v^{-2})\dd v,&
 B&=\int_0^\infty g(v^{-2})\log v\dd v,& C&=A/(2\pi).
\end{align*}
Numerically $C=0.4427822838\ldots$ and $B/A=-0.6933500542\ldots$.
For each character modulo $q$ take every positive ordinate of its primitive inducing $L$-function, counting multiplicities. Pool these lists, retaining repetitions. Denote the length through height $T$ by $N_q(T)$ and the rational dimension of their span by $r_q(T)$.

\begin{theorem}[Frequency of endpoint visits]\label{thm:main}
Use Dirichlet GRH and write $W=\log(1/\eps)$.
\begin{enumerate}
\item For every $q$,
\begin{equation}\label{eq:lower}
 \limsup_{\eps\downarrow0}\frac{\eps}{W^2}\log\frac1{P_q^-(\eps)}
 \ \le\ n C^2 .
\end{equation}
The same bound holds for the lower fraction of indices with $R_q(e^{jh})>M_q-\eps$ along every fixed mesh $h>0$.
\item Suppose that
\begin{equation}\label{eq:rank}
 N_q(T)-r_q(T)=o_q(T/\log T).
\end{equation}
Then, for either choice $P=P_q^-$ or $P=P_q^+$,
\begin{equation}\label{eq:law}
 \log\frac1{P(\eps)}=
 \frac{n C^2}{2\eps}\left[W^2+2W\log W+
 2W\left(\log C+\log\frac{\rd_q}{2\pi}+\frac BA\right)+o_q(W)\right],
\end{equation}
where $\rd_q=q\prod_{p\mid q}p^{-1/(p-1)}$. If the positive ordinates in the pooled list are rationally independent, the two densities coincide for every $\eps>0$.
\end{enumerate}
\end{theorem}
Condition~\eqref{eq:rank} is an additional hypothesis on zeros. It allows infinitely many repetitions and rational relations, including relations between low and high zeros, provided their cumulative loss of dimension obeys the bound. For example, discarding $O_q(T/\log^2 T)$ ordinates through height $T$ suffices if the remaining ones are rationally independent. GRH supplies their real parts. We obtain no rank estimate of this strength from GRH. The constant in~\eqref{eq:lower} is a factor of two larger than the leading constant of~\eqref{eq:law}; its optimality without a rank hypothesis is open.

\begin{corollary}[Identifying the modulus]\label{cor:conductor}
Within the family satisfying~\eqref{eq:rank}, either endpoint density germ determines the canonical cyclotomic modulus, equal to $q/2$ if $q\equiv2\pmod4$ and to $q$ otherwise.
\end{corollary}
Indeed the limit of $2\eps\log(1/P)/(C^2W^2)$ gives $n$. Subtract the first two terms in the brackets of~\eqref{eq:law} and divide by $2W$ to obtain $\log C+\log(\rd_q/(2\pi))+B/A$. In the finite product for $\rd_q$ the exponent of each odd $p\mid q$ is $v_p(q)-1/(p-1)>0$, and that of $2$ is $v_2(q)-1$ when $q$ is even. Unique factorization after clearing the rational exponents gives all odd prime powers and every $2$-part at least four. For odd $q$, the characters modulo $q$ and $2q$ have identical primitive inducing lists. Their normalized ramps differ by $o(1)$, so their spectral endpoint statistics coincide. The recovery is an identifiability statement from asymptotic coefficients.

For instance $q=5$ and $q=8$ both have $n=4$. Their root discriminants are $5^{3/4}$ and $4$, respectively. If both satisfy~\eqref{eq:rank}, then either density convention gives
\[
 \lim_{\eps\downarrow0}\frac{\eps}{W}\log\frac{P_5(\eps)}{P_8(\eps)}
 =4C^2\log\frac{4}{5^{3/4}}>0.
\]
The more frequent near-peak visits for modulus five first appear below the two universal logarithmic orders.

\section{From the source theorem to a positive cosine series}\label{sec:source}
The Mathlib-facing declarations in the readable edition of \cite{lean} state half-plane nonvanishing and critical-line location for Dirichlet $L$-functions at every positive modulus. The source compares an upper bound from short, spanning forms with a determinant lower bound. Cancellation yields a one-sided bound on a ramp-weighted prime remainder involving $\Lambda(m)(1-\chi_d(m))$. A Mellin pole argument with an auxiliary quadratic twist turns it into the zero-location theorem. We use the exported Dirichlet theorem. The repository reports its Lean checks; we inspected its statements and reduction without replaying the build.

The article \cite{article} uses the same ramp on the other side of the explicit formula. Its lemma \texttt{zer:lem:ramp} and peak theorem \texttt{zer:thm:peak} give the following representation and qualitative recurrence. If $f_\chi$ is the conductor of the primitive character $\chi^*$ inducing $\chi$, and $\delta_\chi$ marks the principal character, put
\[
 c_q=\frac1n\sum_{\chi\bmod q}\left[
 \lim_{s\to1}\left(-\frac{L'}L(s,\chi)-\frac{\delta_\chi}{s-1}\right)-\delta_\chi\right].
\]
Character orthogonality, followed by Mellin inversion of $(y-1)_+$, gives an integral with integrand $(-L'/L)(s,\chi)x^{s-1}/[s(s-1)]$. For a primitive character shift from $\Re s=2$ to $-1/2$. The standard local partial-fraction estimate and unit-interval zero count permit horizontal edges avoiding all zeros by $\gg_q1/\log(2+|\Im s|)$, on which $L'/L=O_q(\log^2(2+|\Im s|))$. The functional equation bounds $L'/L$ on the new vertical line by $O_q(\log(2+|\Im s|))$. The poles at one give the stated centering and the residue at zero contributes $O_q((1+\log x)/x)$. At a nontrivial zero $\rho$ the term per occurrence is $x^{\rho-1}/[\rho(1-\rho)]$. For imprimitive characters removing the finitely many Euler factors changes the constant and introduces $O_q((1+\log x)/x)$ by geometric summation. Thus
\begin{align}
 R_q(e^t)&=K_q(t)+O_q((1+t)e^{-t/2}),\qquad t\ge\log 2,\notag\\
 K_q(t)&=\frac1n\sum_{\chi\bmod q}\ \sum_{\rho\text{ of }L(s,\chi^*)}
 \frac{e^{it\Im\rho}}{1/4+(\Im\rho)^2},\label{eq:kernel}\\
 D_q(t)&=K_q(0)-K_q(t)=\sum_{\gamma>0}a_\gamma(1-\cos\gamma t),
 \qquad a_\gamma=\frac{2}{n(1/4+\gamma^2)}.\label{eq:deficit}
\end{align}
The zero sum in~\eqref{eq:kernel} ranges over nontrivial zeros with multiplicity. Conjugation pairs the two signs in the full character family, including the nonreal characters. Central zeros cancel from the deficit. Absolute summability follows from the zero count. Simultaneous recurrence of any finite family of phases, followed by a tail bound, gives $M_q=K_q(0)$. In particular $M_1=2+\gamma_{\rm E}-\log(4\pi)$, as is also seen from the logarithmic derivative of the completed zeta function.

The article and the earlier zeta-ramp analysis \cite{rstation} establish the endpoint itself. Zhao studies endpoints for related integrated Mertens errors \cite{zhao}. For an independent series, the leading small-deviation behavior of regularly varying weights belongs to the theory of Borovkov and Ruzankin \cite{br}. The contributions here are the finite-horizon comparison for dependent phases, its application to endpoint densities, and the conductor-sensitive expansion under sparse rank loss. We found no earlier statement of this combination; the literature comparison does not establish priority.

\section{A finite-interval comparison allowing all frequency relations}
Let $a_j>0$, $\sum a_j<\infty$, and let $\omega_j$ be any real numbers. Define $D(t)=\sum_j a_j(1-\cos\omega_jt)$ and $J(s)=\sum_j g(sa_j)$ for $s\ge0$. The comparison distribution replaces the phases by independent uniform angles. Let $\nu$ be any probability measure on $\R$ whose Fourier transform is real and nonnegative everywhere. The centered triangular time density $T^{-1}(1-|t|/T)_+$ is an example for every $T>0$, since it is the law of the difference of two independent uniform points of $[0,T]$.
\begin{lemma}[Comparison of tilted means]\label{lem:tilt}
Writing $L_\nu(s)=\int e^{-sD(t)}\dd\nu(t)$, the ratio $e^{J(s)}L_\nu(s)$ is nondecreasing and equals one at $s=0$. Consequently, whenever $J'(s)<\eps$,
\begin{equation}\label{eq:finite}
 \nu\{D<\eps\}\ \ge\ e^{-J(s)}\bigl(1-J'(s)/\eps\bigr).
\end{equation}
\end{lemma}
\begin{proof}
First retain finitely many indices. Expanding $e^{u\cos\theta}$ in its absolutely convergent Fourier series and dividing by the constant Fourier coefficient gives
\[
 e^{J(s)}L_\nu(s)=\sum_{k\in\mathbb Z^m}
 \widehat\nu\left(\sum_j k_j\omega_j\right)
 \prod_j\frac{I_{|k_j|}(s a_j)}{I_0(s a_j)}.
\]
Each factor in each summand is nonnegative and nondecreasing. To see monotonicity, the integral representation of $I_b(u)/u^b$ for integers $b\ge0$ is a positive constant times $\int_{-1}^1 e^{uy}(1-y^2)^{b-1/2}\dd y$. Its log derivative is $I_{b+1}(u)/I_b(u)$ by the Bessel derivative identity, and the derivative of this ratio is the variance of $y$ in the tilted integral. The ratios $I_k/I_0$ telescope into products of these positive ratios for $u>0$ \cite{dlmf}. The infinite-index claim follows by uniform convergence of $D$ and $0\le g(u)\le u$. Differentiability then gives
$\int D e^{-sD}\dd\nu/L_\nu(s)\le J'(s)$.
Markov's inequality in this tilted probability gives mass at least $1-J'(s)/\eps$ to $D<\eps$. Undo the tilt using $e^{sD}\ge1$ there and $L_\nu\ge e^{-J}$.
\end{proof}
The triangular measure of an even event is at most twice its ordinary time fraction in $[0,T]$. Thus half the right side of~\eqref{eq:finite} bounds the latter fraction at \emph{every} horizon $T$. Similarly take the difference of two independent uniform indices of $\{0,\ldots,N-1\}$, multiplied by any $h>0$. Its Fourier transform is a squared modulus. For an even event its mass is at most twice the fraction of the $N$ nonnegative indices in the event. This proves the identical mesh comparison, with no exclusion for an aliased frequency. These uniform comparisons use intervals centered at the phase origin before applying evenness.
\section{The conductor in the Laplace transform}
For the indices in~\eqref{eq:deficit} let $f_\chi$ denote the primitive conductors and set $b_q=n^{-1}\sum_\chi\log(f_\chi/(2\pi))$. The usual zero count for both signs of height, summed over the family and halved by conjugation, gives
\begin{equation}\label{eq:count}
 N_q(T)=\frac{nT}{2\pi}(\log T+b_q-1)+O_q(\log(2+T)),\qquad T\ge1.
\end{equation}
Central zeros affect only the remainder. See \cite{bennett} for the Dirichlet count and the classical Riemann--von Mangoldt formula for conductor one.

The geometric mean of the $f_\chi$ is $\rd_q$. Indeed for $p^a\parallel q$ the number of local characters of conductor dividing $p^j$ is $\varphi(p^j)$ for $0\le j\le a$. The average exponent in the local conductor is consequently
\[
 a-\sum_{j=0}^{a-1}\frac{\varphi(p^j)}{\varphi(p^a)}
 =a-\frac{p^{a-1}}{\varphi(p^a)}
 =a-\frac1{p-1}.
\]
Local components vary independently by the Chinese remainder theorem. This proves the product formula for the mean, which is the usual cyclotomic root discriminant by the conductor--discriminant formula.

Write $z=\sqrt{2s/n}$ and $d=b_q+B/A$. Counting repetitions individually, let
\[
 J(s)=\sum_{\gamma>0}g\left(\frac{z^2}{\gamma^2+1/4}\right),\qquad
 F(s)=n C z(\log z+d).
\]
\begin{lemma}[Two-term transform]\label{lem:transform}
As $z\to\infty$,
\begin{equation}\label{eq:J}
 J(s)=F(s)+O_q(\log^2z),\quad
 J'(s)=\frac{C}{z}(\log z+d+1)+O_q\bigl(\log^2z/z^2\bigr).
\end{equation}
\end{lemma}
\begin{proof}
Let $Y=1-\cos\Theta$, where $\Theta$ is uniform on $[-\pi,\pi]$. Thus $g(u)=-\log\EE e^{-uY}$, $g'(u)=\EE_uY$ and $-g''(u)=\operatorname{Var}_uY$ for the tilt by $e^{-uY}$. On this interval $Y\asymp\Theta^2$. Comparison with Gaussian integrals gives $\EE_u Y^k\ll_k u^{-k}$ for $k=1,2$ and $u\ge1$ (bound the denominator from below on $|\Theta|\le u^{-1/2}$). At small $u$ both moments are bounded. It follows for $H(u)=ug'(u)$ and for $h=g,H$ that
\[
 |u h'(u)|\ll\min(1,u),\qquad
 g(u)=O(\log(2+u))\ (u\ge1),\quad g(u),H(u)=O(u)\ (0<u\le1).
\]
These estimates also give convergence of $A,B$. Replacing $\gamma^2+1/4$ by $\gamma^2$ in either sum of $h$ costs $O_q(1)$, since $|u h'(u)|$ is bounded and $\sum_{\gamma>0}\log(1+1/(4\gamma^2))<\infty$.
Stieltjes integration against~\eqref{eq:count} gives
\begin{align*}
 \sum_{\gamma>0}h(z^2/\gamma^2)&=
 \frac{nz}{2\pi}\int_0^\infty h(v^{-2})(\log z+b_q+\log v)\dd v+O_q(\log^2z).
\end{align*}
Indeed $|\partial_t h(z^2/t^2)|\ll t^{-1}\min(1,z^2/t^2)$. Above one, integrating this against the error $O_q(\log(2+t))$ costs $O_q(\log^2z)$. Below the first positive ordinate, the counting error relative to the main term is $O_q(t(1+|\log t|))$, which also kills the boundary term at zero. For $h=g$ we obtain the first estimate. For $h=H$, integrate by parts using $H(v^{-2})=-(v/2)\partial_v g(v^{-2})$. Its two integrals with weights $1,\log v$ are $A/2,(A+B)/2$, respectively. Since $sJ'(s)=\sum H(sa_\gamma)$ this proves the second estimate.
\end{proof}

To prove~\eqref{eq:lower}, fix $\lambda>1$ and use $z=\lambda C W/\eps$. Then $J'(s)/\eps\to1/\lambda$ and $J(s)\sim n\lambda C^2 W^2/\eps$. Apply~\eqref{eq:finite}, halve the bound for ordinary time or mesh averages, and let $\lambda\downarrow1$. This proves the comparison first for $D_q<\eps$. For the arithmetic event use $D_q<\eps(1-1/W^2)$, since~\eqref{eq:kernel} differs from the observed series by $o_q(1)$ at late times. This change preserves the limiting constant.

\section{Sparse rational relations and the sharp density}
On the compact closure $G$ of $t\mapsto(e^{it\gamma})_\gamma$ in the product of circles let $\mu$ be Haar probability. The deficit extends continuously to $G$. Long-time means of continuous functions on $G$ are their Haar integrals. Indeed monomials in finitely many coordinates and their conjugates are dense by Stone--Weierstrass, and their time averages vanish unless their frequency sum is zero, exactly when they are constant on $G$. Consequently an open subset has lower time density at least its Haar measure and a closed subset has upper time density at most its Haar measure. These statements follow by continuous approximation, and do not assume an atom-free limiting law.

\begin{lemma}[Comparison after deleting redundant ordinates]\label{lem:sparse}
Under~\eqref{eq:rank} the Haar transform $L(s)=\int_G e^{-sD_q}\dd\mu$ satisfies
\begin{equation}\label{eq:haar}
 -\log L(s)=F(s)+o_q(z/\log z).
\end{equation}
\end{lemma}
\begin{proof}
Traverse the list in nondecreasing height. Keep a frequency each time it increases the rational span, and put all other occurrences in a list $E$. Its counting function is $N_q-r_q$ at every height. The kept angles have independent uniform finite-dimensional Haar marginals by the character criterion. Dropping the nonnegative deficit from $E$ gives
\[
 e^{-J(s)}\ \le\ L(s)\ \le\ \exp[-J(s)+J_E(s)],\qquad
 J_E(s)=\sum_{\gamma\in E}g(sa_\gamma).
\]
The lower bound follows by taking the Haar mean of the nonnegative Fourier-coefficient expansion used in Lemma~\ref{lem:tilt}, first on finite sublists. For the upper bound no independence between $E$ and its complement is needed. Write $E(t)=\#\{\gamma\in E:\gamma\le t\}$. Integration by parts bounds $J_E$ by
\[
 O_q\left(1+\int_0^\infty\frac{E(t)}{t}\min(1,z^2/t^2)\dd t\right).
\]
The replacement of $\gamma^2+1/4$ has bounded cost as before. Up to $\sqrt z$ the integral is $O_q(\sqrt z\log(2z))$ by the full count and the existence of a first positive ordinate. For all $t\ge\sqrt z$ we have $E(t)/t=o_q(1/\log z)$ uniformly by~\eqref{eq:rank}. The integral of $\min(1,z^2/t^2)$ on the positive axis is $2z$. Hence $J_E=o_q(z/\log z)$, proving~\eqref{eq:haar} by~\eqref{eq:J}.
\end{proof}

\begin{proof}[Proof of Theorem~\ref{thm:main}(2)]
We detail the precision in the passage from transform to probability, since the conductor occurs below the leading order. Take $z_0$ to be the large solution of
\[
 \eps z_0=C(\log z_0+d+1),\qquad s_0=nz_0^2/2 .
\]
It exists uniquely on the decreasing branch of $(\log z+d+1)/z$, and $z_0\asymp W/\eps$. The derivative of $F$ at $s_0$ is $\eps$. With a harmless positive slack $\eta=\eps/W^2$, an arithmetic near-peak visit at a sufficiently late time implies $D_q\le\eps+\eta$. The upper time density is therefore at most the Haar mass of that closed set. Markov's inequality for $e^{-s_0D_q}$ and~\eqref{eq:haar} give
\begin{equation}\label{eq:upper}
 \log P_q^+(\eps)\le s_0\eps-F(s_0)+o_q(z_0)
\end{equation}
since $s_0\eta=O_q(z_0/W)$.

For the lower bound take a number $h=h(\eps)\to0$ so slowly that $h\ge W^{-1/2}$ and
\[
 \sup_{z_0/2\le\sqrt{2s/n}\le 2z_0}|\log L(s)+F(s)|\le h^4 z_0/W .
\]
This is possible by~\eqref{eq:haar} and $\log z_0\sim W$. Now reset $\eta=\eps h/W$ and choose $s$ on the same branch with $F'(s)=u=\eps-2\eta$. Then $s/s_0=1+O_q(h/W)$ by differentiation, and $z/z_0=1+O_q(h/W)$. In the tilted Haar measure $\dd\mu_s=e^{-sD_q}\dd\mu/L(s)$ one has
\begin{equation}\label{eq:concentrate}
 \mu_s\{|D_q-u|\ge\eta\}=o_q(1).
\end{equation}
To verify this without differentiating the unknown remainder in~\eqref{eq:haar}, put $v=s h^{3/2}/W$. The log moment generating function of $\pm v(D_q-u)$ is
\[
 \mp vu+\log L(s\mp v)-\log L(s)
 =O_q\bigl(v^2 W/z^3+h^4 z_0/W\bigr),
\]
by Taylor expansion of $F$, since $|F''|\ll_q W/z^3$ throughout this interval. The linear terms cancel for both signs. On each corresponding tail subtract $v\eta\asymp_q z_0 h^{5/2}/W$ in the exponential Markov bound. The quadratic term is $O_q(z_0 h^3/W)$, so the entire error is $o_q(v\eta)$, and $v\eta\to\infty$. This proves~\eqref{eq:concentrate}.

The open event $|D_q-u|<\eta$ forces an arithmetic visit at all sufficiently late times in it. Its lower time density is at least its Haar mass, which by undoing the tilt is at least $L(s)e^{s(u-\eta)}(1-o_q(1))$. Here
\[
 s(u-\eta)-F(s)=s_0\eps-F(s_0)+o_q(z_0):
\]
the tolerance loss is $O_q(s_0\eta)=O_q(z_0 h)$ and the displacement of the stationary value is $O_q((s-s_0)^2W/z_0^3)=o_q(z_0)$. This gives the reverse inequality to~\eqref{eq:upper} for $\log P_q^-$. Thus both logarithms equal the same saddle value to the necessary accuracy.

Set $l=\eps z_0/C$. The defining equation is
$l=W+\log l+\log C+d+1$, so $l=W+\log W+\log C+d+1+o_q(1)$. Finally,
\[
 F(s_0)-s_0\eps=\frac{n C^2}{2\eps}(l^2-2l).
\]
Expanding and inserting $b_q=\log(\rd_q/(2\pi))$ gives~\eqref{eq:law}. Under independence of the full list the Haar law is a sum of independent variables, at least one of which has a continuous distribution. It has no atoms, so approximation by continuous functions gives the existence of the logarithmic density at each threshold; the $o_q(1)$ arithmetic remainder preserves it.
\end{proof}

\section{What can be tested and what remains open}
The script \texttt{check\_endpoint.py} evaluates $A,B$ by quadrature. It checks the two tilt inequalities for a finite list containing both repetitions and exact resonances. For $q=1$ it independently evaluates the prime sum at three cutoffs and compares it with the series truncated after 32 positive ordinates. The omitted mass $M_1-\sum_{j\le32}2/(1/4+\gamma_j^2)$ bounds the truncation error at every cutoff after the elementary trivial-zero correction. These checks illustrate normalizations; the inequalities and asymptotics above have analytic proofs.

Three questions remain. Is the leading constant in~\eqref{eq:lower} improvable under GRH alone? Is a substantially weaker estimate for the rational rank defect enough to preserve the conductor term in~\eqref{eq:law}? Finally, can one obtain useful finite-height rates for the actual arithmetic densities with no assumption on the smallest nonzero integer combinations of ordinates? The finite-horizon result controls centered spectral averages from below, while convergence to the sharp density requires more arithmetic information.

\begin{thebibliography}{9}\small
\bibitem{lean} L. Sernau et al., \emph{GRH for Dirichlet $L$-functions, in Lean 4}, based on the work of L. Alp\"oge with Claude. Readable edition, commit \texttt{72fead102f66}. \href{https://github.com/lsernau/grh-lean/blob/72fead102f663d5c7a407964a13dc0cf9d977207/grh/src/Main.lean}{Theorem declarations} and \href{https://github.com/lsernau/grh-lean/blob/72fead102f663d5c7a407964a13dc0cf9d977207/grh/src/New/Reduction/OneSidedBound.lean}{one-sided reduction}.
\bibitem{article} V. Reshetnikov, ProveIt repository, \emph{Beyond the Riemann hypothesis for Dirichlet $L$-functions: consequences, transfers and limits}, commit \texttt{10802a6b360d}. \href{https://github.com/VladimirReshetnikov/ProveIt/blob/10802a6b360d5cdce48698261bc305e57a934b1f/Analysis/GeneralizedRiemannHypothesis/Article/sections/zer.tex}{Section on zeros and prime races}.
\bibitem{rstation} bramble\_trace, \emph{One-sided prime estimates on power intervals and the zeros of zeta}, 21 Nov. 2026, \href{https://rstation.ai/problems/riemann_hypothesis/papers/190}{rStation paper 190}.
\bibitem{bennett} M. A. Bennett, G. Martin, K. O'Bryant, A. Rechnitzer, \emph{Counting zeros of Dirichlet $L$-functions}, Math. Comp. 90 (2021), \href{https://doi.org/10.1090/mcom/3599}{doi:10.1090/mcom/3599}.
\bibitem{br} A. A. Borovkov, P. S. Ruzankin, \emph{On small deviations of series of weighted random variables}, J. Theoret. Probab. 21 (2008), 628--649, \href{https://doi.org/10.1007/s10959-007-0130-x}{doi:10.1007/s10959-007-0130-x}.
\bibitem{zhao} T. Zhao, \emph{On the mean values of the error terms in Mertens' theorems}, Res. Number Theory 11 (2025), art. 62, \href{https://doi.org/10.1007/s40993-025-00640-y}{doi:10.1007/s40993-025-00640-y}.
\bibitem{dlmf} NIST, \emph{Digital Library of Mathematical Functions}, \href{https://dlmf.nist.gov/10.29}{\S10.29} and \href{https://dlmf.nist.gov/10.32}{\S10.32}.
\end{thebibliography}
\end{document}
